\documentclass[11pt,a4paper]{article}
\usepackage[utf8]{inputenc}
\usepackage[T1]{fontenc}
\usepackage{lmodern}
\usepackage{amsmath,amssymb}
\usepackage{booktabs}
\usepackage{graphicx}
\usepackage{microtype}
\usepackage{hyperref}
\usepackage{url}
\usepackage{enumitem}

\hypersetup{
    colorlinks=true,
    linkcolor=blue,
    citecolor=blue,
    urlcolor=blue,
    pdftitle={SCM: Sleep-Consolidated Memory for Language Agents},
    pdfauthor={SCM Research Team}
}

\title{SCM: Sleep-Consolidated Memory for Language Agents\\
\large A Reproducible Evaluation of One-Shot Encoding, Sleep-Stage Consolidation, and Selective Forgetting}
\author{SCM Research Team}
\date{May 1, 2026}

\begin{document}
\maketitle

\begin{abstract}
Large language models remain constrained by memory systems that are either bounded by context windows or implemented as continuously growing retrieval stores without lifecycle management. We present SCM (Sleep-Consolidated Memory), a memory architecture that combines bounded working memory, event-aware encoding, staged sleep consolidation, selective forgetting, and contradiction-safe version lineage. We evaluate SCM using repository-backed benchmark artifacts and reproducibility packs. In the latest human-memory benchmark, SCM attains one-shot recall accuracy 1.0, micro-sleep disambiguation gain 1.0, deep-sleep disambiguation gain 1.0, deep-sleep noise retention 0.0, and contradiction-versioning accuracy 1.0. In a 10-seed baseline stress comparison with 96 memory pairs per seed, lexical retrieval, vector retrieval, and SCM without sleep remain at disambiguation recall 0.0, while SCM with deep sleep reaches mean disambiguation recall 0.9052 with mean noise retention 0.0. In long-horizon evaluation, sleep-enabled SCM reaches final family-aware disambiguation recall 0.9812 and final noise retention 0.0, versus awake-only recall 0.0 and noise retention 1.0. Reproducibility-pack reruns pass end-to-end. These results support SCM as a strong and testable memory-lifecycle architecture and establish a rigorous baseline for external user studies.
\end{abstract}

\section{Introduction}
Long-context prompting and retrieval-augmented memory have improved practical agent performance, but they do not fully address long-horizon memory quality under conflict and noise. Context-window approaches degrade when relevant information is not near prompt boundaries \cite{liu2023lost}. Retrieval stores improve recall but are commonly append-only and weak on consolidation and forgetting \cite{lewis2020retrieval}. Hierarchical memory systems such as MemGPT improve scale control but are not explicitly modeled as biologically inspired memory lifecycles \cite{packer2023memgpt}. Production memory layers such as Mem0 improve personalization, but still operate as always-awake retrieval systems \cite{chhikara2025mem0}.

Human memory is dynamic and lifecycle-based. It combines bounded working capacity \cite{miller1956magical}, sleep-based consolidation \cite{rasch2013sleep}, and active forgetting to preserve signal quality \cite{tononi2003sleep}. SCM adopts these principles as an engineering architecture for language agents: memory is encoded, replayed, consolidated, selectively forgotten, and contradiction-updated with lineage preservation.

This paper asks a narrow question: does a sleep-centered memory lifecycle produce measurable gains over awake-only and standard retrieval controls under stress testing? We answer using current SCM artifacts and show strong gains in disambiguation and selective noise suppression while preserving one-shot and contradiction-update behaviors.

\section{Contributions}
This paper makes four concrete contributions:
\begin{enumerate}[leftmargin=1.5em]
    \item It defines SCM as an end-to-end memory lifecycle with explicit micro-sleep and deep-sleep stages, rather than an always-awake retrieval wrapper.
    \item It introduces contradiction-safe version lineage where conflicting updates supersede older traces while preserving auditability.
    \item It provides controlled stress evidence showing that sleep stages, not baseline lexical or vector retrieval mechanics, drive disambiguation gains.
    \item It ships a reproducibility pathway that replays baseline comparison, long-horizon behavior, guardrails, and smoke stability checks from recorded artifacts.
\end{enumerate}

\section{Method}
\subsection{Architecture Overview}
SCM combines bounded working memory with persistent graph-based long-term memory. Incoming dialogue is transformed into typed concepts with contextual tags, then written to working memory and long-term memory with value signals. Retrieval uses semantic cues and spreading activation. Sleep orchestration periodically runs consolidation and forgetting passes, then syncs the updated memory state back into long-term storage.

The architecture is organized into six project phases:
\begin{itemize}[leftmargin=1.5em]
    \item \textbf{Phase 1}: selective encoding and grasp estimation.
    \item \textbf{Phase 2}: event binding and association growth.
    \item \textbf{Phase 3}: spreading-activation retrieval and hypothesis ranking.
    \item \textbf{Phase 4}: micro-sleep and deep-sleep consolidation.
    \item \textbf{Phase 5}: forgetting dynamics and contradiction-safe versioning.
    \item \textbf{Phase 6}: evaluation harnesses, guardrails, and reproducibility pack.
\end{itemize}

\subsection{Sleep Stages}
Micro-sleep performs lightweight replay, local reinforcement, and duplicate-pressure cleanup. It targets near-term disambiguation improvements with low disruption.

Deep sleep performs broader replay and consolidation, relation synthesis, and forgetting dynamics. Retired concept states are preserved in sync payloads to maintain contradiction-version lineage through consolidation.

\subsection{Selective Forgetting and Version Lineage}
Forgetting dynamics compute retention from grasp, salience, rehearsal, association density, recency, and interference. Concepts transition through \texttt{ACTIVE}, \texttt{SUPPRESSED}, and \texttt{ARCHIVED}.

Contradictions are handled through version lineage instead of overwrite semantics: a new conflicting trace can supersede prior versions while preserving root and parent links for traceability.

\section{Experimental Protocol}
All values reported in this paper are extracted directly from benchmark artifacts:
\begin{itemize}[leftmargin=1.5em]
    \item Human-memory suite: \url{research/metrics/phase6_human_memory_latest.json}
    \item Baseline stress comparison: \url{research/metrics/baseline_comparison_latest.json}
    \item Long-horizon behavior: \url{research/metrics/long_horizon_latest.json}
    \item Guardrails: \url{research/metrics/phase6_guardrails_latest.json}
    \item Backend smoke: \url{research/metrics/phase6_backend_smoke_latest.json}
    \item Reproducibility pack: \url{research/reproducibility/reproducibility_pack_latest.json}
\end{itemize}

We report metric values and pass flags exactly as stored by these harnesses.

\section{Results}
\subsection{Human-Memory Suite}
The latest Phase 6 human-memory artifact reports full pass status. Core metrics are shown in Table~\ref{tab:human_suite}.

\begin{table}[htbp]
\centering
\caption{Phase 6 human-memory benchmark metrics.}
\label{tab:human_suite}
\begin{tabular}{@{}lc@{}}
\toprule
\textbf{Metric} & \textbf{Value} \\
\midrule
One-shot recall accuracy & 1.0000 \\
Micro-sleep disambiguation gain (absolute) & 1.0000 \\
Deep-sleep disambiguation gain (absolute) & 1.0000 \\
Deep-sleep noise retention & 0.0000 \\
Contradiction-versioning accuracy & 1.0000 \\
Overall pass & True \\
\bottomrule
\end{tabular}
\end{table}

These results indicate that SCM simultaneously preserves one-shot encoding, post-sleep disambiguation, and contradiction-safe update behavior.

\subsection{Baseline Stress Comparison}
Table~\ref{tab:baseline} summarizes the 10-seed stress comparison with 96 memory pairs per seed. Lexical control, vector control, and SCM awake baseline all stay at disambiguation recall 0.0. Sleep-enabled variants provide large gains.

\begin{table}[htbp]
\centering
\caption{Behavioral stress comparison (10 seeds, 96 pairs per seed).}
\label{tab:baseline}
\begin{tabular}{@{}lc@{}}
\toprule
\textbf{Setting} & \textbf{Disambiguation Recall Mean} \\
\midrule
Lexical retrieval baseline & 0.0000 \\
Vector retrieval baseline & 0.0000 \\
SCM awake-only baseline & 0.0000 \\
SCM micro-sleep & 1.0000 \\
SCM deep-sleep & 0.9052 \\
\midrule
SCM deep-sleep noise retention mean & 0.0000 \\
SCM deep-sleep pass rate & 1.0000 \\
\bottomrule
\end{tabular}
\end{table}

The attribution pattern is clear: large gains appear only when sleep consolidation stages are enabled.

\subsection{Long-Horizon Behavior}
In long-horizon evaluation, the final-day comparison between sleep-enabled and awake-only modes is shown in Table~\ref{tab:long_horizon}.

\begin{table}[htbp]
\centering
\caption{Final long-horizon comparison (day 7).}
\label{tab:long_horizon}
\begin{tabular}{@{}lcc@{}}
\toprule
\textbf{Metric} & \textbf{Awake-Only} & \textbf{Sleep-Enabled} \\
\midrule
Family-aware disambiguation recall & 0.0000 & 0.9812 \\
Noise retention & 1.0000 & 0.0000 \\
Anchor hit rate & 1.0000 & 1.0000 \\
\midrule
Disambiguation lift (sleep - awake) & \multicolumn{2}{c}{0.9812} \\
Noise reduction (awake - sleep) & \multicolumn{2}{c}{1.0000} \\
Strict duplicate-pair lift & \multicolumn{2}{c}{0.0000} \\
\bottomrule
\end{tabular}
\end{table}

The strict duplicate-pair lift remains a transparent stress limitation and is not hidden in reporting.

\subsection{Reliability and Guardrails}
The guardrail suite reports \texttt{overall\_pass=True}, with zero hits for tracked deprecation patterns (\texttt{datetime.utcnow(} and \texttt{.dict(}). The backend smoke artifact reports 20/20 HTTP success, 20/20 overall pass, average latency 45.2 ms, p50 latency 22.7 ms, and p90 latency 26.3 ms.

The reproducibility pack reports \texttt{overall\_pass=True}, with passing summaries across baseline comparison, long-horizon behavior, guardrails, and smoke pytests.

\section{Discussion}
Across suites, SCM is strongest where memory quality depends on lifecycle transformations, especially disambiguation under duplicate pressure and selective noise suppression over time. Because awake-only and no-sleep controls remain at or near zero disambiguation in stress settings, the observed gains are unlikely to be explained by retrieval plumbing alone.

The current evidence is internal but rigorous: deterministic harnesses, multi-seed stress tests, regression guardrails, and reproducibility-pack reruns. This supports SCM as an engineering-science result with a clear path to broader external validation.

\section{Limitations}
This work has four primary limitations.
\begin{enumerate}[leftmargin=1.5em]
    \item Evaluations are mostly synthetic and harness-driven rather than broad naturalistic user studies.
    \item External comparisons remain limited; wider cross-system benchmarks should be added.
    \item Strict duplicate-pair behavior under long-horizon pressure remains weaker than family-aware behavior.
    \item The paper does not claim consciousness or full human cognition; it evaluates bounded human-like memory dynamics only.
\end{enumerate}

\section{Conclusion}
SCM demonstrates that a sleep-consolidated memory lifecycle can produce large, reproducible gains over awake-only and standard retrieval controls on targeted memory tasks. The architecture preserves one-shot behavior and contradiction-safe updates while improving long-horizon disambiguation and removing low-value noise. The next step is external user validation followed by broader baseline expansion.

\section*{Reproducibility Statement}
The benchmark and regression scripts are available in the repository, including:
\begin{itemize}[leftmargin=1.5em]
    \item \url{tests/baseline_comparison.py}
    \item \url{tests/long_horizon_benchmark.py}
    \item \url{tests/phase6_guardrails.py}
    \item \url{tests/reproducibility_pack.py}
\end{itemize}
The latest reproducibility summary is stored at \url{research/reproducibility/reproducibility_pack_latest.json}, which reports \texttt{overall\_pass=True} as of May 1, 2026.

\bibliographystyle{plain}
\bibliography{references}

\end{document}
